DIGITAL FILTERS - LECTURE 6

FILTERS AND DIFFERENCE EQUATIONS

$$y[n]=\sum_{k=0}^{M}b_kx[n-k]-\sum_{k=1}^{N}a_ky[n-k]$$

$$H(z)=\frac{\sum_{k=0}^{M}b_kz^{-k}}{1+\sum_{k=1}^{N}a_kz^{-k}}$$

$$H(e^{j\omega})=H(z)|_{z=e^{j\omega}}$$

Frequency response gives steady-state output of a sinusoid:

$$x[n]=A\cos(\omega_0n+\phi)$$

$$y[n]=A|H(e^{j\omega_0})|$$

$$\cos(\omega_0n+\phi+\theta(\omega_0))$$

Here $\theta(\omega)=\arg H(e^{j\omega})$.

Gain in dB:

$$G(\omega)=20\log_{10}|H(e^{j\omega})|$$

Equivalent:

$$G(\omega)=10\log_{10}|H(e^{j\omega})|^2$$

At the -3 dB cutoff for unity passband gain:

$$|H(e^{j\omega_c})|=\frac{1}{\sqrt{2}}$$

$$|H(e^{j\omega_c})|^2=\frac12$$

$\omega=2\pi f/F_s$ radians/sample.

IDEAL FILTER FREQUENCY RANGES

Lowpass: pass $|\omega|<\omega_c$; stop above $\omega_c$ to $\pi$.

Highpass: stop below $\omega_c$; pass above $\omega_c$ to $\pi$.

Bandpass: pass $\omega_1<|\omega|<\omega_2$.

Bandstop: reject $\omega_1<|\omega|<\omega_2$.

PHASE DELAY AND GROUP DELAY

$$\tau_p(\omega)=-\frac{\theta(\omega)}{\omega}$$

$$\tau_g(\omega)=-\frac{d\theta(\omega)}{d\omega}$$

$\tau_p$ is phase delay, $\tau_g$ is envelope/group delay, measured in samples. Use unwrapped phase when differentiating.

Pure delay:

$$H(e^{j\omega})=e^{-j\omega D}$$

$$\theta(\omega)=-\omega D$$

$$\tau_p=\tau_g=D$$

FIR FILTER AND LINEAR PHASE

$$H(z)=\sum_{n=0}^{N}h[n]z^{-n}$$

There are N+1 coefficients and linear-phase delay $D=N/2$ for a symmetric or antisymmetric FIR (away from phase jumps).

Symmetric:

$$h[n]=h[N-n]$$

Antisymmetric:

$$h[n]=-h[N-n]$$

Linear phase factor:

$$H(e^{j\omega})=e^{-j\omega N/2}A(\omega)$$

For symmetric FIR $A(\omega)$ is real. For antisymmetric FIR the amplitude factor is imaginary.

FOUR LINEAR-PHASE FIR TYPES

Type I: symmetric, odd number of taps (N even). No forced zero at DC or Nyquist.

Type II: symmetric, even number of taps (N odd). Must have H(-1)=0, i.e., zero at $\omega=\pi$.

Type III: antisymmetric, odd number of taps (N even). Zeros at $\omega=0$ and $\omega=\pi$.

Type IV: antisymmetric, even number of taps (N odd). Must have zero at $\omega=0$.

SIMPLE FIRST-ORDER FIR LOWPASS

$$H_{LP}(z)=\frac12(1+z^{-1})$$

$$H_{LP}(e^{j\omega})=e^{-j\omega/2}\cos(\omega/2)$$

For $0\le\omega\le\pi$:

$$|H_{LP}|=\cos(\omega/2)$$

$$\omega_c=\frac{\pi}{2}$$

M lowpass sections in cascade:

$$H_M(z)=[H_{LP}(z)]^M$$

$$|H_M|=\cos^M(\omega/2)$$

$$\omega_c=2\arccos(2^{-1/(2M)})$$

SIMPLE FIRST-ORDER FIR HIGHPASS

$$H_{HP}(z)=\frac12(1-z^{-1})$$

$$H_{HP}(e^{j\omega})=je^{-j\omega/2}\sin(\omega/2)$$

$$|H_{HP}|=\sin(\omega/2)$$

FIRST-ORDER IIR LOWPASS

$$H_{LP}(z)=\frac{1-\alpha}{2}\frac{1+z^{-1}}{1-\alpha z^{-1}}$$

$$|H_{LP}|^2=\frac{(1-\alpha)^2(1+\cos\omega)}{2(1+\alpha^2-2\alpha\cos\omega)}$$

For the -3 dB cutoff:

$$\cos\omega_c=\frac{2\alpha}{1+\alpha^2}$$

$$\alpha=\frac{\cos\omega_c}{1+\sin\omega_c}$$

Stable causal case: $|\alpha|<1$.

FIRST-ORDER IIR HIGHPASS

$$H_{HP}(z)=\frac{1+\alpha}{2}\frac{1-z^{-1}}{1-\alpha z^{-1}}$$

$$|H_{HP}|^2=\frac{(1+\alpha)^2(1-\cos\omega)}{2(1+\alpha^2-2\alpha\cos\omega)}$$

Same -3 dB design relation:

$$\cos\omega_c=\frac{2\alpha}{1+\alpha^2}$$

SECOND-ORDER IIR BANDPASS

$$H_{BP}(z)=\frac{1-\alpha}{2}$$

$$\frac{1-z^{-2}}{1-\beta(1+\alpha)z^{-1}+\alpha z^{-2}}$$

Center frequency:

$$\omega_0=\arccos(\beta)$$

Bandwidth between its -3 dB frequencies:

$$B_\omega=\arccos\left(\frac{2\alpha}{1+\alpha^2}\right)$$

SECOND-ORDER IIR BANDSTOP

$$H_{BS}(z)=\frac{1+\alpha}{2}$$

$$\frac{1-2\beta z^{-1}+z^{-2}}{1-\beta(1+\alpha)z^{-1}+\alpha z^{-2}}$$

Notch center:

$$\omega_0=\arccos(\beta)$$

POLE-ZERO INTERPRETATION

$$H(z)=C\frac{\prod_{m=1}^{M}(z-\zeta_m)}{\prod_{k=1}^{N}(z-p_k)}z^{N-M}$$

At $z=e^{j\omega}$:

$$|H|=|C|\frac{\prod_m|e^{j\omega}-\zeta_m|}{\prod_k|e^{j\omega}-p_k|}$$

A zero on the unit circle blocks its frequency. A pole near the unit circle tends to strengthen nearby frequencies. For causal stability, all poles must be inside $|z|<1$.
